\[
R_l^m(\mathbf r_i)
=\frac{r_i^l}{(l+m)!}
P_l^m(\cos\theta_i)e^{im\phi_i}
\]

we already allocated space for real and imaginary arrays, so we fill them like this

\begin{itemize}
    \item Diagonal terms
    \item Off-Diagonal terms
    \item General terms
\end{itemize}

\subsubsection*{The Diagonal Terms (m = l)}

For the diagonal terms, we use the standard property of Associated Legendre Polynomials

\[
P_l^l(\cos\theta)
=-(2l-1)\sin\theta\,P_{l-1}^{l-1}(\cos\theta)
\]

\[
\begin{aligned}
R_l^l
&=\frac{r^l}{(2l)!}
\left[-(2l-1)\sin\theta\,P_{l-1}^{l-1}(\cos\theta)\right]\\
&\qquad\times e^{il\phi}
\end{aligned}
\]

\[
\begin{aligned}
R_l^l
&=-\frac{r\sin\theta\,e^{i\phi}}{2l}\\
&\qquad\times\left[\frac{r^{l-1}}{(2l-1)!}
P_{l-1}^{l-1}(\cos\theta)e^{i(l-1)\phi}\right]
\end{aligned}
\]

\[
R_l^l=-\frac{x+iy}{2l}R_{l-1}^{l-1}
\]

\subsubsection*{The Off-Diagonal Terms (m = l - 1)}

For the first sub-diagonal, the Legendre polynomial property is:

\[
P_l^{l-1}(\cos\theta)
=(2l-1)\cos\theta\,P_{l-1}^{l-1}(\cos\theta)
\]

\[
\begin{aligned}
R_l^{l-1}
&=\frac{r^l}{(2l-1)!}
\left[(2l-1)\cos\theta\,P_{l-1}^{l-1}(\cos\theta)\right]\\
&\qquad\times e^{i(l-1)\phi}
\end{aligned}
\]

\[
\begin{aligned}
R_l^{l-1}
&=r\cos\theta\\
&\qquad\times\left[\frac{r^{l-1}}{(2l-1)!}
P_{l-1}^{l-1}(\cos\theta)e^{i(l-1)\phi}\right]
\end{aligned}
\]

\[
R_l^{l-1}=zR_{l-1}^{l-1}
\]

\subsubsection{The General Terms (0 $\leq$ m $\leq$ l - 2)}

For all remaining terms, we rely on Bonnet's recursion formula for Associated Legendre Polynomials:

\[
\begin{aligned}
(l-m)P_l^m(\cos\theta)
&=(2l-1)\cos\theta\,P_{l-1}^m(\cos\theta)\\
&\qquad-(l+m-1)P_{l-2}^m(\cos\theta)
\end{aligned}
\]

\[
\begin{aligned}
\frac{(l-m)r^lP_l^m e^{im\phi}}{(l+m)!}
&=\frac{(2l-1)zr^{l-1}P_{l-1}^m e^{im\phi}}{(l+m)!}\\
&\qquad-\frac{(l+m-1)r^lP_{l-2}^m e^{im\phi}}{(l+m)!}
\end{aligned}
\]

\[
\begin{aligned}
\frac{(2l-1)z}{l+m}
&\left[\frac{r^{l-1}P_{l-1}^m e^{im\phi}}{(l+m-1)!}\right]\\
&=\frac{(2l-1)z}{l+m}R_{l-1}^m
\end{aligned}
\]

\[
\begin{aligned}
\frac{r^2(l+m-1)}{(l+m)(l+m-1)(l+m-2)!}
&\left[r^{l-2}P_{l-2}^m e^{im\phi}\right]\\
&=\frac{r^2}{l+m}R_{l-2}^m
\end{aligned}
\]

\[
\begin{aligned}
(l-m)R_l^m
&=\frac{(2l-1)zR_{l-1}^m-r^2R_{l-2}^m}{l+m}
\end{aligned}
\]

\[
\begin{aligned}
R_l^m
&=\frac{(2l-1)zR_{l-1}^m-r^2R_{l-2}^m}{(l-m)(l+m)}
\end{aligned}
\]

Now we have everything for building Multipole coefficients.
